\section{おわりに}
\label{sec_conclusion}
本稿では，大喜利のお題画像における文字の扱いの変遷と，その影響をまとめた．
読める文字の有無を表す文字ありは，改訂のたびに観測できる基準に近づき，適合率0.94，再現率0.99で判定できるようになった．
一方，題として扱う文字を表すお題文は，投稿者の意図に基づく定義を保ったため，教師と生徒の誤りが文字の出所の判断に集中した．
そこで，文字を写っている場所による4つの層と，文の形による5つの形に分け，お題文ありを属性から規則で導く定義を提案した．
提案の属性は同じ判断者の付け直しで記号の94\%が一致し，既定の規則 $ r_{1b} $ で導いたお題文ありは第5版と全データの推定で98.5\%一致した．
また，規則の選択だけでお題文ありの割合が7.40\%から16.67\%まで変わることを示した．
一方，Qwen3-VL-32B は層を区別できず，1,620枚で学習した生徒も既存のパッケージに届かなかった．
今後は，層を生徒に任せて教師には文の形だけを問う構成で教師ラベルを増やすこと，人どうしの一致を測ること，ボケの内容から規則の妥当性を確かめることが課題である．
